% =========================
% report/discussion.tex
% =========================

\subsection{جمع‌بندی علمی از یافته‌ها}
یافته‌ها باید بر اساس شواهد کیفی (هم‌پوشانی‌ها) و شواهد کمی (جدول معیارها، توزیع خطا، تحلیل زمان) تفسیر شوند. نسبت درون‌نهشتی بالا معمولاً نشانهٔ سازگاری هندسی است، اما تضمین قطعی «درستی» نیست؛ زیرا ممکن است تطبیق‌های غلط به‌صورت منسجم نیز یک هموگرافی غلط بسازند. بنابراین، ترکیب چند شاخص (InlierRatio + تعداد تطبیق + خطای بازفرافکنی در صورت وجود زمین‌حقیقت) برای قضاوت قابل‌اعتمادتر است.

\subsection{محدودیت‌ها}
\begin{itemize}
  \item \textbf{وابستگی به بافت/ساختار:} روش‌های ویژگی‌محور در سطوح بسیار یکنواخت عملکرد ضعیف‌تری دارند.
  \item \textbf{حساسیت به بازتاب و نور شدید:} تغییرات شدید نور می‌تواند نقاط کلیدی را ناپایدار کند.
  \item \textbf{نقض فرض صفحه‌ای:} در سطوح سه‌بعدی پیچیده، هموگرافی یکپارچه کافی نیست.
  \item \textbf{تکرار الگوها:} الگوهای تکراری باعث تطبیق‌های غلط منسجم می‌شوند.
\end{itemize}

\subsection{کارهای آینده}
\begin{itemize}
  \item افزودن تخمین «کیفیت/اعتماد» صریح و آستانه‌گذاری برای ارجاع به بازبینی انسانی،
  \item استفاده از روش‌های ترکیبی (مثلاً ویژگی + قیود هندسی اضافی یا شناسایی ROI)،
  \item در صورت نیاز به دقت بالاتر، بررسی روش‌های یادگیری (با دادهٔ کافی و زمین‌حقیقت).
\end{itemize}
